% Pista ana-25s-q4a · Anàlisi
% Procedència: PAU Catalunya 2025 · Convocatòria de setembre · Sèrie 3 · Exercici 4 — Opció A
\documentclass[11pt,a4paper]{article}
\usepackage{pista}
\pistainfo{Catalunya 2025}{Convocatòria de setembre}{Sèrie 3}

\begin{document}

\section*{\textcolor{blauPista}{Exercici 4 --- Opció A}}

\noindent Un vidrier està reparant un dels vitralls de la Sagrada Família la forma del qual és la de la part ombrejada de la figura adjunta. S'ha adonat que Gaudí el va dissenyar de manera que un dels costats segueix la funció $y = 3\sin(x/4)$ i un altre segueix la funció $y = 3\cos(x/4)$, on $x$ i $y$ estan expressades en metres.

\subsection*{\textit{a)} \normalfont Raoneu a quina funció correspon cada gràfica i calculeu les coordenades dels punts $B$ i $C$ assenyalats a la figura (tenint en compte que $A$ és l'origen de coordenades). \hfill\textnormal{[1 punt]}}

\pista{Avalua les dues funcions a $x = 0$: una d'elles passa per l'origen $A$ i l'altra per $B$, això t'indica directament quina correspon a cada costat del vitrall. El punt $C$ és la intersecció de les dues corbes: iguala $3\sin(x/4) = 3\cos(x/4)$ per trobar la seva coordenada $x$.}

\subsection*{\textit{b)} \normalfont Calculeu el preu del vitrall sabent que costa 750\,\euro/m$^{2}$. \hfill\textnormal{[1,5 punts]}}

\pista{L'àrea d'una regió entre dues corbes és la integral definida de la diferència de les dues funcions (la de ``dalt'' menys la de ``baix''). Però has de trobar quins són els extrems de la integral definida: són les coordenades $x$ dels punts $A$ i $C$ de l'apartat anterior. Com sempre, has de rellegir la pregunta, perquè no en tens prou amb calcular l'àrea. Un cop calculada, l'has de multiplicar pel preu del metre quadrat i ara sí que obtindràs la resposta.}

\end{document}
